%% chapters/ch4_results.tex

\chapter{實驗結果與討論}
\label{ch:results}

\section{結果}
\label{sec:results}

這一節呈現實驗結果。若圖的邊緣有多餘留白，可用 \verb|\figtr| 裁切，
第五個參數是 \verb|左 下 右 上| 四個值，未標單位時以大點（bp）計：

\figtr{0.6}{figures/example_fig}{裁切後的示範圖。}{fig:trimmed}{0 0 0 0}

圖 \ref{fig:trimmed} 示範裁切參數的寫法。

\section{討論}
\label{sec:discussion}

這一節討論結果的意義與限制，並與既有工作 \cite{reckdahl2006graphics}
及第 \ref{ch:background} 章的理論基礎相互對照。

程式碼片段可用 \verb|lstlisting| 環境，裡面的內容不算進 thesis-cli stats 的字數：

\begin{lstlisting}
thesis-cli build --strict
thesis-cli check-refs --strict
\end{lstlisting}
